%============================================================
% HALAMAN KATA PENGANTAR
%============================================================
\cleardoublepage
\chapter*{KATA PENGANTAR}
\addcontentsline{toc}{chapter}{KATA PENGANTAR}
\thispagestyle{fancy}

Puji syukur penulis panjatkan ke hadirat Tuhan Yang Maha Esa karena atas rahmat dan penyertaan-Nya penulis dapat menyelesaikan buku Proyek Akhir yang berjudul ``Fraud Detection dalam Data Perawatan Kesehatan melalui Metode Machine Learning'' dengan baik.

Penyusunan buku Proyek Akhir ini dilakukan sebagai salah satu syarat untuk menyelesaikan pendidikan pada Program Studi Sarjana Terapan Teknik Informatika, Politeknik Elektronika Negeri Surabaya. Penelitian ini membahas pendekatan machine learning berbasis claim-centric untuk mendeteksi potensi fraud atau upcoding pada klaim JKN/BPJS dalam sistem INA-CBG, khususnya ketika label fraud aktual pada level klaim tidak tersedia.

Dalam proses penyusunan Proyek Akhir ini, penulis memperoleh banyak bantuan, arahan, dan dukungan dari berbagai pihak. Oleh karena itu, penulis menyampaikan terima kasih kepada:

\begin{enumerate}[label=\arabic*., leftmargin=*]
  \item Ibu Rosiyah Faradisa, S.Si., M.Si. selaku dosen pembimbing yang telah memberikan arahan, masukan, dan bimbingan selama proses penelitian dan penulisan.
  \item Ibu Tessy Badriyah, S.Kom., M.T., Ph.D. selaku dosen pembimbing yang telah memberikan dukungan ilmiah, koreksi, dan saran yang sangat berarti dalam penyempurnaan penelitian ini.
  \item Keluarga penulis yang selalu memberikan doa, dukungan, dan semangat selama proses pengerjaan Proyek Akhir.
  \item Teman-teman serta seluruh pihak yang telah membantu penulis, baik secara langsung maupun tidak langsung, dalam menyelesaikan penelitian ini.
\end{enumerate}

Penulis menyadari bahwa buku Proyek Akhir ini masih memiliki keterbatasan. Oleh karena itu, penulis sangat terbuka terhadap kritik dan saran yang membangun untuk pengembangan penelitian ini di masa mendatang.

Akhir kata, penulis berharap semoga buku Proyek Akhir ini dapat memberikan manfaat bagi pengembangan analitik data klaim kesehatan, khususnya sebagai dasar prioritisasi verifikasi klaim berisiko tinggi.

\vspace{1cm}
\begin{flushright}
Surabaya, 25 April 2026\\[1.5cm]
Penulis
\end{flushright}
